\begin{trapbox}
\begin{description}[style=nextline, leftmargin=2.8em, labelsep=0.5em, itemsep=0.35em]
    \item[\textbf{\textcolor{CTrap}{易错点一（忽略取等）}}] 误以为不等式在任何 $x$ 处等号都可能成立。
    \item[\textbf{\textcolor{CTrap}{正确理解（唯一零点）：}}] 等号仅在 $x=0$ 时成立，$x>0$ 时均为严格不等式。
    \item[\textbf{\textcolor{CTrap}{错因分析（条件遗漏）：}}] 未注意导数零点唯一性及函数单调性，导致等号范围扩大。
\end{description}

\medskip
\begin{description}[style=nextline, leftmargin=2.8em, labelsep=0.5em, itemsep=0.35em]
    \item[\textbf{\textcolor{CTrap}{易错点二（范围滥用）}}] 将不等式随意用于 $x<0$ 的情况。
    \item[\textbf{\textcolor{CTrap}{正确理解（非负定义域）：}}] 所有不等式的前提是 $x\geq0$，$x<0$ 时结论可能不成立。
    \item[\textbf{\textcolor{CTrap}{错因分析（前提忽视）：}}] 未审清题目中变量的范围，盲目套用公式。
\end{description}

\medskip
\begin{description}[style=nextline, leftmargin=2.8em, labelsep=0.5em, itemsep=0.35em]
    \item[\textbf{\textcolor{CTrap}{易错点三（精度混淆）}}] 记忆混乱，放缩时选错上界或下界，导致方向错误或精度不足。
    \item[\textbf{\textcolor{CTrap}{正确理解（按需选取）：}}] 根据题目要求（需上界还是下界、精度高低）选择合适的不等式。
    \item[\textbf{\textcolor{CTrap}{错因分析（概念模糊）：}}] 未理解各不等式的特点及相互强弱关系。
\end{description}
\end{trapbox}
